\chapter{Source Code Excerpts}
\label{appendix:code}

This appendix presents condensed excerpts of the framework's core implementation. The complete source code is available in the accompanying repository. Excerpts are edited for brevity: logging statements, input validation, and purely cosmetic code are omitted where they do not affect functional understanding.

% ------------------------------------------------------------------
\section{Domain-Agnostic Simulation Engine}
\label{sec:code-engine}

Listing~\ref{lst:app-base-federate} shows the complete \texttt{BaseFederate} abstract class, which defines the three-method contract that all domain modules must implement. The engine manages the HELICS lifecycle, time advancement, and report formatting independently of domain logic.

\begin{listing}[H]
\begin{minted}[fontsize=\scriptsize]{python}
import helics as h
from abc import ABC, abstractmethod
from typing import Dict

class BaseFederate(ABC):
    """Abstract base for all domain federates.
    Subclasses must implement:
        initialize_components()
        run_simulation()
        generate_report() -> Dict
    """

    def __init__(self, name: str, use_helics: bool = False,
                 time_step: float = 1.0, sim_duration: float = 3600.0):
        self.name = name
        self.federate = None
        self.use_helics = use_helics
        self.time_step = time_step
        self.sim_duration = sim_duration

    def setup_federate(self):
        """Setup HELICS federate (no-op in standalone mode)."""
        if not self.use_helics:
            return
        fedinfo = h.helicsCreateFederateInfo()
        h.helicsFederateInfoSetCoreName(fedinfo, self.name)
        h.helicsFederateInfoSetCoreTypeFromString(fedinfo, "zmq")
        h.helicsFederateInfoSetCoreInitString(
            fedinfo, "--federates=1 --autobroker")
        h.helicsFederateInfoSetTimeProperty(
            fedinfo, h.helics_property_time_delta, self.time_step)
        self.federate = h.helicsCreateValueFederate(self.name, fedinfo)
        self._register_publications()
        h.helicsFederateEnterExecutingMode(self.federate)

    def _register_publications(self):
        """Override to register domain-specific HELICS publications."""
        pass

    def advance_time(self, current_time: float) -> float:
        """Advance simulation clock by one time step."""
        if self.use_helics and self.federate:
            return h.helicsFederateRequestTime(
                self.federate, current_time + self.time_step)
        return current_time + self.time_step

    def cleanup(self):
        """Cleanup HELICS federate."""
        if self.use_helics and self.federate:
            h.helicsFederateFinalize(self.federate)
            h.helicsFederateFree(self.federate)
            h.helicsCloseLibrary()

    # --- Contract methods (domain modules implement these) ---

    @abstractmethod
    def initialize_components(self):
        """Create domain-specific agents and models."""
        ...

    @abstractmethod
    def run_simulation(self):
        """Execute the main simulation loop."""
        ...

    @abstractmethod
    def generate_report(self) -> Dict:
        """Return dict with 'scenario', 'components', 'metrics' keys."""
        ...
\end{minted}
\caption{\texttt{engine/base.py} --- Domain-agnostic simulation engine (complete)}
\label{lst:app-base-federate}
\end{listing}

% ------------------------------------------------------------------
\section{Domain Module: Smart Grid (E1)}
\label{sec:code-e1}

Listing~\ref{lst:app-e1-federate} shows the condensed \texttt{SmartGridFederate} class, demonstrating how a domain module implements the engine contract. Agent models (solar PV, wind turbine, battery) and the main simulation loop illustrate the pattern followed by all domain modules.

\begin{listing}[H]
\begin{minted}[fontsize=\scriptsize]{python}
import numpy as np
from engine.base import BaseFederate

class SolarPV:
    def __init__(self, pv_id, capacity_kw, bus_id):
        self.capacity_kw = capacity_kw
        self.generation_kw = 0.0

    def calculate_generation(self, hour, cloud_factor=1.0):
        if 6 <= hour <= 18:
            time_factor = np.sin((hour - 6) * np.pi / 12)
            self.generation_kw = self.capacity_kw * time_factor * cloud_factor * 0.85
        else:
            self.generation_kw = 0.0
        return self.generation_kw

class WindTurbine:
    def __init__(self, turbine_id, capacity_kw, bus_id):
        self.capacity_kw = capacity_kw

    def calculate_generation(self, wind_speed_ms):
        cut_in, rated, cut_out = 3.0, 12.0, 25.0
        if wind_speed_ms < cut_in or wind_speed_ms > cut_out:
            return 0.0
        elif wind_speed_ms < rated:
            return self.capacity_kw * ((wind_speed_ms - cut_in) / (rated - cut_in)) ** 3
        return self.capacity_kw

class BatteryStorage:
    def __init__(self, battery_id, capacity_kwh, power_kw, bus_id):
        self.capacity_kwh = capacity_kwh
        self.power_kw = power_kw
        self.soc = 0.5
        self.soc_min, self.soc_max = 0.2, 0.95
        self.efficiency = 0.92

    def charge(self, power_kw, dt_hours):
        max_energy = min(power_kw, self.power_kw) * dt_hours * self.efficiency
        available = (self.soc_max - self.soc) * self.capacity_kwh
        energy = min(max_energy, available)
        self.soc += energy / self.capacity_kwh
        return energy / dt_hours if dt_hours > 0 else 0

    def discharge(self, power_kw, dt_hours):
        max_energy = min(power_kw, self.power_kw) * dt_hours
        available = (self.soc - self.soc_min) * self.capacity_kwh
        energy = min(max_energy, available) * self.efficiency
        self.soc -= energy / (self.capacity_kwh * self.efficiency)
        return energy / dt_hours if dt_hours > 0 else 0

class SmartGridFederate(BaseFederate):
    def __init__(self, name, use_helics=False, num_solar_pvs=50,
                 num_wind_turbines=3, num_batteries=5, num_loads=800, **kwargs):
        super().__init__(name, use_helics, time_step=900.0, sim_duration=24*3600)
        self.num_solar_pvs = num_solar_pvs
        self.num_wind_turbines = num_wind_turbines
        self.num_batteries = num_batteries
        self.solar_pvs, self.wind_turbines, self.batteries = [], [], []

    def initialize_components(self):
        for i in range(self.num_solar_pvs):
            capacity = np.random.uniform(5.0, 10.0)
            self.solar_pvs.append(SolarPV(i, capacity, bus_id=i % 33))
        for i in range(self.num_wind_turbines):
            self.wind_turbines.append(WindTurbine(i, 500.0, bus_id=i))
        for i in range(self.num_batteries):
            self.batteries.append(BatteryStorage(i, 50.0, 25.0, i))

    def run_simulation(self):
        current_time = 0.0
        while current_time < self.sim_duration:
            hour = (current_time / 3600.0) % 24
            cloud_factor = 0.8 + 0.2 * np.random.random()
            solar_total = sum(pv.calculate_generation(hour, cloud_factor)
                              for pv in self.solar_pvs)
            wind_speed = 5.0 + 7.0 * np.random.random()
            wind_total = sum(wt.calculate_generation(wind_speed)
                             for wt in self.wind_turbines)
            total_renewable = solar_total + wind_total
            # Load calculation, battery dispatch, voltage profile ...
            current_time = self.advance_time(current_time)

    def generate_report(self):
        return {
            "scenario": "E1 - Smart Grid with Renewable Integration",
            "components": {"solar_pv_count": len(self.solar_pvs),
                           "wind_turbine_count": len(self.wind_turbines),
                           "battery_count": len(self.batteries)},
            "metrics": {"renewable_penetration_pct": ...,
                        "total_curtailment_kwh": ...}
        }
\end{minted}
\caption{\texttt{scenarios/scenario\_e1.py} --- Smart grid federate (condensed)}
\label{lst:app-e1-federate}
\end{listing}

% ------------------------------------------------------------------
\section{Strategy Pattern: Charging Strategies (E2)}
\label{sec:code-strategy}

Listing~\ref{lst:app-strategy} illustrates the strategy pattern implementation in Scenario~E2. The \texttt{ChargingStrategy} enum selects between behavioural variants without modifying the simulation engine or the charging infrastructure model.

\begin{listing}[H]
\begin{minted}[fontsize=\scriptsize]{python}
from enum import Enum
from engine.base import BaseFederate

class ChargingStrategy(Enum):
    UNCOORDINATED = "uncoordinated"  # Immediate charging on arrival
    SMART = "smart"                  # Grid-aware coordinated charging
    V2G = "v2g"                      # Vehicle-to-Grid bidirectional

class ElectricVehicle:
    def __init__(self, ev_id, battery_capacity_kwh, initial_soc):
        self.battery_capacity_kwh = battery_capacity_kwh
        self.soc = initial_soc
        self.soc_target = 0.8
        self.v2g_enabled = np.random.random() > 0.5

    def charge(self, power_kw, dt_hours, price_per_kwh):
        available = (0.9 - self.soc) * self.battery_capacity_kwh
        energy = min(power_kw * dt_hours * 0.90, available)
        self.soc += energy / self.battery_capacity_kwh
        return energy / dt_hours

    def discharge(self, power_kw, dt_hours, price_per_kwh):
        """V2G discharge (only if v2g_enabled)."""
        if not self.v2g_enabled:
            return 0.0
        available = (self.soc - 0.3) * self.battery_capacity_kwh
        energy = min(power_kw * dt_hours, available) * 0.90
        self.soc -= energy / (self.battery_capacity_kwh * 0.90)
        return energy / dt_hours

class SmartChargingController:
    """Grid-aware charging coordinator."""
    def allocate(self, evs, grid_capacity, base_load, hour):
        available_power = grid_capacity - base_load
        queue = sorted(evs, key=lambda ev: (ev.soc, ev.departure_time))
        for ev in queue:
            alloc = min(ev.max_charge_rate_kw, available_power / max(1, len(queue)))
            yield ev, alloc

class EVChargingFederate(BaseFederate):
    def __init__(self, name, use_helics=False,
                 strategy=ChargingStrategy.SMART, **kwargs):
        super().__init__(name, use_helics, time_step=300.0, sim_duration=24*3600)
        self.strategy = strategy

    def run_simulation(self):
        current_time = 0.0
        while current_time < self.sim_duration:
            hour = (current_time / 3600.0) % 24
            if self.strategy == ChargingStrategy.UNCOORDINATED:
                self._uncoordinated_step(hour)
            elif self.strategy == ChargingStrategy.SMART:
                self._smart_step(hour)
            else:  # V2G
                self._v2g_step(hour)
            current_time = self.advance_time(current_time)
\end{minted}
\caption{\texttt{scenarios/scenario\_e2.py} --- Strategy pattern for EV charging (condensed)}
\label{lst:app-strategy}
\end{listing}

% ------------------------------------------------------------------
\section{Cross-Domain Coupling: Energy--Mobility (E2+M1)}
\label{sec:code-crossdomain}

Listing~\ref{lst:app-crossdomain} shows the cross-domain scenario architecture. An EV inherits mobility behaviour from the M1 traffic model while carrying an E2 battery model, demonstrating how domain modules interact through shared agent state and module reuse via Python imports.

\begin{listing}[H]
\begin{minted}[fontsize=\scriptsize]{python}
from enum import Enum
from engine.base import BaseFederate
from scenarios.scenario_m1 import (
    TrafficNetwork, Vehicle as M1Vehicle,
    create_random_vehicles, update_signals, move_vehicles_step,
)
from scenarios.scenario_e2 import (
    ElectricVehicle, ChargingStation, SmartChargingController,
    ChargingStrategy,
)

class EVTrafficVehicle(M1Vehicle):
    """EV that drives through traffic to reach a charging station."""

    class Phase(Enum):
        DRIVING_TO_STATION = "driving_to_station"
        CHARGING = "charging"
        DRIVING_BACK = "driving_back"
        DONE = "done"

    def __init__(self, vehicle_id, origin, station_pos, ev, return_dest):
        super().__init__(vehicle_id, origin, station_pos)
        self.ev = ev           # E2 battery model (composition)
        self.station_pos = station_pos
        self.return_destination = return_dest
        self.phase = self.Phase.DRIVING_TO_STATION

    @property
    def at_station(self):
        return self.completed and self.phase == self.Phase.DRIVING_TO_STATION

class CrossDomainE2M1(BaseFederate):
    """Coupled Energy-Mobility simulation on a shared 5x5 traffic grid."""

    def __init__(self, name="CrossE2M1", coupled=True,
                 num_evs=500, num_background=2000, grid_size=5, **kwargs):
        super().__init__(name, time_step=1.0,
                         sim_duration=kwargs.get('sim_duration_hours', 3) * 3600)
        self.coupled = coupled
        self.network = TrafficNetwork(grid_size=grid_size, spacing_m=1250.0)

    def initialize_components(self):
        self._place_charging_stations()
        self._create_evs()
        self._create_background_vehicles()
        if self.coupled:
            self.all_vehicles = self.ev_vehicles + self.background_vehicles
        else:
            self.all_vehicles = self.background_vehicles  # EVs teleport

    def run_simulation(self):
        current_time = 0.0
        while current_time < self.sim_duration:
            # Step 1: Update traffic signals (M1 domain)
            update_signals(self.network.signals, self.all_vehicles, 1.0)
            # Step 2: Move vehicles through network (M1 mechanics)
            move_vehicles_step(self.all_vehicles, self.network, 1.0)
            # Step 3: Handle EV charging transitions (E2 domain)
            for ev_veh in self.ev_vehicles:
                if ev_veh.at_station:
                    self._charge_ev(ev_veh, current_time)
            current_time = self.advance_time(current_time)

    def generate_report(self):
        return {
            "scenario": f"E2+M1 Cross-Domain (coupled={self.coupled})",
            "components": {"num_evs": len(self.ev_vehicles),
                           "num_background": len(self.background_vehicles),
                           "charging_stations": len(self.charging_stations)},
            "metrics": {"avg_travel_time_to_station_s": ...,
                        "co2_emissions_kg": ...,
                        "charging_target_achievement_pct": ...,
                        "peak_grid_load_kw": ...}
        }
\end{minted}
\caption{\texttt{scenarios/scenario\_e2m1.py} --- Cross-domain energy--mobility coupling (condensed)}
\label{lst:app-crossdomain}
\end{listing}

% ------------------------------------------------------------------
\section{Scenario Execution Pattern}
\label{sec:code-runner}

Listing~\ref{lst:app-runner} shows the standard entry-point pattern used by all scenarios. The engine contract ensures that every scenario follows the same lifecycle: initialise, setup federate, simulate, report, cleanup.

\begin{listing}[H]
\begin{minted}[fontsize=\scriptsize]{python}
def run_scenario_e1(use_helics=False, **kwargs):
    """Run Scenario E1 with optional HELICS co-simulation."""
    grid = SmartGridFederate("SmartGrid_E1", use_helics=use_helics, **kwargs)
    grid.initialize_components()
    grid.setup_federate()
    grid.run_simulation()
    report = grid.generate_report()
    grid.cleanup()
    return report

def compare_strategies():
    """Compare multiple strategies using the same engine contract."""
    results = {}
    for strategy in ChargingStrategy:
        fed = EVChargingFederate("EV_E2", strategy=strategy)
        fed.initialize_components()
        fed.setup_federate()
        fed.run_simulation()
        results[strategy.value] = fed.generate_report()
        fed.cleanup()
    return results
\end{minted}
\caption{Standard scenario execution and strategy comparison pattern}
\label{lst:app-runner}
\end{listing}
